%% ============================================================
%% shared/services/compute/ec2.tex
%% Amazon EC2 — Elastic Compute Cloud
%% Relevant to: SAA (core) + SAP (advanced placement / migration)
%% ============================================================

\servicesection{Amazon EC2}{\bothtag}{%
  Virtual servers in the cloud. The most fundamental AWS compute service.}

\subsection{What is EC2?}

EC2 provides resizable compute capacity in the cloud. You rent virtual machines
(called \textbf{instances}) rather than buying physical hardware. You choose the
OS, CPU, RAM, storage, and network configuration.

\begin{keyfacts}
\begin{itemize}
  \item Instances run inside a \textbf{VPC} in a specific \textbf{Availability Zone}
  \item Billing: per-second (Linux) or per-hour (Windows) — minimum 60 seconds
  \item An \textbf{AMI} (Amazon Machine Image) defines the OS + pre-installed software
  \item \textbf{User data} scripts run once at first boot (bootstrap automation)
  \item \textbf{Instance metadata} is accessible at \texttt{169.254.169.254}
\end{itemize}
\end{keyfacts}

% ─────────────────────────────────────────
\subsection{Instance Types}

Instance types follow a naming pattern: \texttt{<family><generation>.<size>}
e.g.\ \texttt{m6i.large}, \texttt{c5.xlarge}, \texttt{r6g.2xlarge}.

\begin{comparebox}[title={Instance family cheat-sheet}]
\begin{tabular}{@{}lll@{}}
\toprule
\textbf{Family} & \textbf{Use case} & \textbf{Examples} \\
\midrule
General Purpose   & Balanced CPU/RAM/network  & \texttt{m6i, t3, t4g} \\
Compute Optimised & High CPU ratio           & \texttt{c6i, c7g}     \\
Memory Optimised  & Large in-memory workloads & \texttt{r6i, x2idn, u-*} \\
Storage Optimised & High sequential I/O       & \texttt{i3, i4i, d3}  \\
Accelerated       & GPU / ML inference        & \texttt{p4, g5, inf2} \\
\bottomrule
\end{tabular}
\end{comparebox}

\begin{examtip}
\textbf{T-family (burstable):} T3/T4g instances accumulate CPU credits when idle
and spend them during spikes. Use \texttt{unlimited} mode to avoid throttling
at extra cost. Exam favourite for variable-load web apps with a low baseline.
\end{examtip}

% ─────────────────────────────────────────
\subsection{Purchasing Options}

\begin{comparebox}[title={Pricing models}]
\begin{tabular}{@{}p{3.2cm}p{5.5cm}p{4cm}@{}}
\toprule
\textbf{Model} & \textbf{Best for} & \textbf{Discount vs On-Demand} \\
\midrule
On-Demand          & Short-term, unpredictable          & —                \\
Reserved (1 or 3yr)& Steady-state, predictable load     & up to 72\%       \\
Savings Plans      & Flexible RI alternative            & up to 66\%       \\
Spot               & Fault-tolerant, flexible workloads & up to 90\%       \\
Dedicated Host     & License compliance (BYOL)          & varies           \\
Dedicated Instance & Single-tenant hardware (no BYOL)   & varies           \\
Capacity Reservation & Reserve AZ capacity (no discount) & —               \\
\bottomrule
\end{tabular}
\end{comparebox}

\begin{examtip}
\textbf{Spot Instance interruption} = 2-minute warning via instance metadata
and EventBridge. Architect around this with \emph{Spot + On-Demand mixed}
Auto Scaling groups. Keywords: "lowest cost," "fault-tolerant," "batch jobs."
\end{examtip}

\begin{gotcha}
Reserved Instances apply to a \textbf{specific AZ} (Zonal RI → capacity reserved)
or \textbf{any AZ in a region} (Regional RI → no capacity reservation but flexible).
Exam will test whether you need the capacity guarantee.
\end{gotcha}

% ─────────────────────────────────────────
\subsection{Storage for EC2}

\begin{keyfacts}
\begin{itemize}
  \item \textbf{EBS (Elastic Block Store)} — network-attached, persists after
        instance stop/terminate (if configured). One EBS $\leftrightarrow$ one AZ.
  \item \textbf{Instance Store} — physically attached NVMe; \emph{ephemeral}
        (lost on stop/terminate). Fastest possible IOPS. Good for caches/scratch.
  \item \textbf{EFS} — shared NFS mount across multiple instances/AZs.
  \item \textbf{FSx for Windows} — SMB share for Windows workloads.
\end{itemize}
\end{keyfacts}

\subsubsection{EBS Volume Types}
\begin{comparebox}[title={EBS types}]
\begin{tabular}{@{}llll@{}}
\toprule
\textbf{Type}  & \textbf{Use case}          & \textbf{Max IOPS} & \textbf{Max throughput} \\
\midrule
gp3 (SSD)      & General purpose boot/dev   & 16,000            & 1,000 MB/s \\
gp2 (SSD)      & Legacy general purpose     & 16,000            & 250 MB/s   \\
io2 Block Expr & Critical DBs (sub-ms)      & 256,000           & 4,000 MB/s \\
io1 (SSD)      & I/O intensive DBs          & 64,000            & 1,000 MB/s \\
st1 (HDD)      & Throughput-heavy big data  & 500               & 500 MB/s   \\
sc1 (HDD)      & Cold, infrequent access    & 250               & 250 MB/s   \\
\bottomrule
\end{tabular}
\end{comparebox}

\begin{examtip}
\textbf{gp3 vs gp2:} gp3 lets you provision IOPS and throughput \emph{independently}.
gp2 ties IOPS to size (3 IOPS/GB, max 16,000). Default to gp3 for new workloads.
\end{examtip}

% ─────────────────────────────────────────
\subsection{Networking}

\begin{keyfacts}
\begin{itemize}
  \item Each instance gets a \textbf{primary private IP} in its subnet (persists)
  \item \textbf{Elastic IP (EIP)} — static public IP you own; reassignable.
        \cost{Charged when \emph{not} attached to a running instance.}
  \item \textbf{ENI (Elastic Network Interface)} — virtual NIC; you can attach
        multiple ENIs to an instance for network separation.
  \item \textbf{Enhanced Networking} — SR-IOV for higher throughput, lower latency.
        Enabled via ENA driver (most modern instances) or Intel 82599 VF.
  \item \textbf{Placement Groups}:
    \begin{itemize}
      \item \emph{Cluster} — same rack, lowest latency, 10 Gbps between instances.
      \item \emph{Spread} — different hardware, max 7 instances/AZ.
      \item \emph{Partition} — racks isolated per partition; for Hadoop/Kafka/HDFS.
    \end{itemize}
\end{itemize}
\end{keyfacts}

% ─────────────────────────────────────────
\subsection{EC2 Auto Scaling}

Auto Scaling maintains the \emph{desired capacity}, replacing unhealthy instances
and scaling in/out based on policies.

\begin{keyfacts}
\begin{itemize}
  \item \textbf{Launch Template} (preferred) or Launch Configuration (legacy)
  \item \textbf{Scaling policies}:
    \begin{itemize}
      \item \emph{Target Tracking} — keep metric at a value (e.g.\ CPU = 50\%).
            Simplest; AWS manages scale-in/out automatically.
      \item \emph{Step Scaling} — scale by different amounts at different thresholds.
      \item \emph{Scheduled} — predictable load patterns (e.g.\ weekday 9am).
      \item \emph{Predictive} — ML-based, pre-warms capacity.
    \end{itemize}
  \item \textbf{Cooldown period} — default 300 s; prevents thrashing.
  \item \textbf{Warm-up period} — time before new instance counts toward metrics.
  \item \textbf{Lifecycle hooks} — pause instance at launch/terminate for custom logic.
\end{itemize}
\end{keyfacts}

\begin{examtip}
  \textbf{"Scale fast, scale slow":} Target tracking with a low scale-out cooldown
  handles sudden spikes. Add a longer scale-in cooldown to avoid premature
  termination. Pair with an \emph{ALB} for HTTP(S) and \emph{NLB} for TCP.
\end{examtip}

% ─────────────────────────────────────────
\subsection{Security}

\begin{keyfacts}
\begin{itemize}
  \item \textbf{Security Groups} — stateful firewall at the instance/ENI level.
        Default: allow all outbound, deny all inbound.
  \item \textbf{IAM Instance Profile} — grants the instance an IAM role.
        \emph{Never} store credentials on an instance; use roles instead.
  \item \textbf{Key Pairs} — RSA key for SSH (Linux) / RDP password decrypt (Windows).
        AWS stores the public key only.
  \item \textbf{EC2 Instance Connect} — browser-based SSH without long-lived keys.
  \item \textbf{SSM Session Manager} — shell access without SSH, no open port 22
        needed. Preferred for secure environments.
\end{itemize}
\end{keyfacts}

\begin{saadepth}
For the SAA exam, understand the \textbf{Security Group vs NACL} distinction:
Security Groups are stateful (return traffic auto-allowed), NACLs are stateless
(return traffic needs explicit allow). NACLs apply at the subnet level.
\end{saadepth}

\begin{sapdepth}
SAP scenarios involve \textbf{migrating large EC2 fleets}:
\begin{itemize}
  \item \awsservice{AWS Migration Hub} — track migration progress centrally.
  \item \awsservice{AWS Application Migration Service (MGN)} — lift-and-shift at scale;
        replicates disks continuously, minimal downtime cutover.
  \item \awsservice{VM Import/Export} — one-off migrations from VMware/Hyper-V.
  \item Right-sizing: use \awsservice{AWS Compute Optimizer} (ML recommendations)
        or \awsservice{Trusted Advisor} for over-provisioned instances.
  \item \textbf{Mixed-instance policies} in ASG: combine Spot + On-Demand with
        multiple instance families/sizes for cost and availability.
\end{itemize}
\end{sapdepth}

\begin{gotcha}
Stopping vs Terminating: a \textbf{stopped} instance retains its EBS root volume;
a \textbf{terminated} instance deletes it (unless \texttt{DeleteOnTermination=false}).
Hibernate preserves RAM to EBS — instance must be EBS-backed, encryption enabled.
\end{gotcha}
